\begin{trapbox}
\begin{description}[style=nextline, leftmargin=2.8em, labelsep=0.5em, itemsep=0.35em]
    \item[\textbf{\textcolor{CTrap}{易错点一（底数范围混淆）}}] 忽视底数 $a$ 必须满足 $a>0$ 且 $a\neq 1$，误以为对任意实数 $a$ 成立。
    \item[\textbf{\textcolor{CTrap}{正确理解（正底非一）：}}] 结论仅在 $a>0$ 且 $a\neq 1$ 时成立。$a=1$ 时函数为常函数，不等式变为等式但无研究价值；$a\leq 0$ 时指数函数无定义或性质不同。
    \item[\textbf{\textcolor{CTrap}{错因分析（定义域遗忘）：}}] 忽略指数函数的基本定义域和底数限制，机械记忆公式导致错误。
\end{description}

\medskip
\begin{description}[style=nextline, leftmargin=2.8em, labelsep=0.5em, itemsep=0.35em]
    \item[\textbf{\textcolor{CTrap}{易错点二（等号条件遗漏）}}] 应用不等式后忽略讨论等号成立条件，或错误认为等号总成立。
    \item[\textbf{\textcolor{CTrap}{正确理解（中点重合）：}}] 等号成立当且仅当 $x=y$，即两点横坐标相同（函数图像上两点重合）。
    \item[\textbf{\textcolor{CTrap}{错因分析（凸性理解不足）：}}] 对“严格下凸”理解不深，误以为在某些情况下 $x\neq y$ 也可能取等。
\end{description}

\medskip
\begin{description}[style=nextline, leftmargin=2.8em, labelsep=0.5em, itemsep=0.35em]
    \item[\textbf{\textcolor{CTrap}{易错点三（方向记反）}}] 记错不等式方向，误写为 $\frac{a^{x}+a^{y}}{2} \leq a^{\frac{x+y}{2}}$。
    \item[\textbf{\textcolor{CTrap}{正确理解（算术平均大）：}}] 由于函数下凸，两端点算术平均大于中点函数值，故不等号为 $\geq$。
    \item[\textbf{\textcolor{CTrap}{错因分析（图像想象错误）：}}] 混淆“上凸”与“下凸”的几何直观，或将指数函数误认为上凸函数。
\end{description}
\end{trapbox}
